\chapter{Cholesterol Levels}
\section*{What Is This Test?}
The cholesterol test, commonly part of a fasting lipid panel, measures total cholesterol, low-density lipoprotein cholesterol (LDL-C), high-density lipoprotein cholesterol (HDL-C), and triglycerides. Total cholesterol = LDL-C + HDL-C + VLDL-C (where VLDL-C is estimated as triglycerides/5 when TG <400 mg/dL). LDL-C is calculated using the Friedewald equation: LDL-C = Total cholesterol --- HDL-C --- (Triglycerides/5) \cite{friedewald1972estimation}. This equation is invalid when triglycerides >400 mg/dL (direct LDL measurement or the Martin-Hopkins equation is used) \cite{martin2013comparison}. LDL-C is the primary atherogenic lipoprotein and the primary target of lipid-lowering therapy. Statin therapy intensity (high, moderate, low) is based on the degree of LDL-C lowering required. Per the 2018 ACC/AHA multisociety guideline on the management of blood cholesterol \cite{grundy2019guideline}, the following groups are candidates for statin therapy: (1) clinical atherosclerotic cardiovascular disease (ASCVD) --- high-intensity statin, (2) LDL-C $\geq$190 mg/dL --- high-intensity statin, (3) diabetes (age 40--75) with LDL-C 70--189 mg/dL --- moderate- to high-intensity statin, and (4) primary prevention (age 40--75, LDL-C 70--189 mg/dL, 10-year ASCVD risk $\geq$7.5\%) --- moderate- to high-intensity statin. HDL-C is inversely associated with ASCVD risk; HDL <40 mg/dL in men and <50 mg/dL in women is a risk factor.
\section*{Alternative Names}
Lipid Panel; Lipid Profile; Fasting Lipid Profile; Total Cholesterol; LDL; HDL; Triglycerides
\section*{Principle}
Total cholesterol and triglycerides are measured by enzymatic colorimetric methods (cholesterol oxidase and glycerol kinase/lipase). HDL-C is measured by a direct homogeneous assay (polyanion/detergent selective solubilization) or after precipitation of apoB-containing lipoproteins. LDL-C is usually calculated (Friedewald equation, if TG <400) or measured by direct homogeneous assay.
\section*{History}
The Framingham Heart Study, reporting from the late 1950s onward, established blood cholesterol as a major risk factor for coronary heart disease \cite{wilson1998prediction}. In 1972, Friedewald and colleagues introduced an equation that allowed LDL-C to be estimated from total cholesterol, HDL-C, and triglycerides without ultracentrifugation. The 1984 Lipid Research Clinics Coronary Primary Prevention Trial and the 1985 NIH Consensus Conference established cholesterol lowering as a therapeutic goal, and the first statin, lovastatin, was approved in 1987. The Scandinavian Simvastatin Survival Study (4S, 1994) demonstrated that statin therapy reduced mortality in secondary prevention, reshaping guideline strategy. The 2013 ACC/AHA guideline shifted from fixed LDL targets to statin intensity based on calculated risk, and the 2018 ACC/AHA multisociety guideline refined this framework, incorporating non-HDL cholesterol and newer evidence.
\section*{Why This Test Is Ordered}
\begin{itemize}
  \item Routine cardiovascular risk assessment and screening in adults, as part of the lipid panel
  \item Estimation of 10-year ASCVD risk with the Pooled Cohort Equations to guide statin decisions
  \item Evaluation of patients with established atherosclerotic cardiovascular disease (ASCVD) to guide secondary prevention
  \item Screening for familial hypercholesterolemia when LDL-C is markedly elevated (typically $\geq$190 mg/dL)
  \item Detection and monitoring of hypertriglyceridemia, including assessment of pancreatitis risk when triglycerides are very high
  \item Monitoring response to statin therapy and other lipid-lowering agents
  \item Assessment of cardiovascular risk in patients with diabetes, hypertension, or other major risk factors
\end{itemize}
\section*{Primary vs Secondary Test}
The lipid panel is a primary screening test for cardiovascular risk assessment in adults and is routinely obtained in primary prevention. When results are abnormal, secondary and specialty tests such as apolipoprotein B (apoB), lipoprotein(a), and coronary artery calcium scoring further refine risk and guide therapy.
\section*{How This Test Is Performed}
Blood is collected by venipuncture into a serum tube or a plasma tube containing EDTA or heparin after a 9--12 hour fast, which is preferred for a complete panel including triglycerides. Total cholesterol and triglycerides are measured by enzymatic colorimetric methods (cholesterol oxidase and glycerol kinase), and HDL-C by a direct homogeneous assay. LDL-C is calculated with the Friedewald equation when triglycerides are below 400 mg/dL or measured directly by homogeneous assay otherwise. Results are typically available within 1 day.
\section*{What Preparation Is Needed}
A fast of 9--12 hours (water only) is required for a full lipid panel, especially when triglycerides are to be interpreted. The patient should maintain a usual diet for the 2 weeks before testing, avoid alcohol for 24 hours, and refrain from strenuous exercise immediately before sampling. Testing should be postponed during acute illness, because intercurrent disease transiently lowers cholesterol. Non-fasting sampling is acceptable for initial screening in some current guidelines.
\section*{Contraindications and Precautions}
\begin{itemize}
  \item The Friedewald equation is invalid when triglycerides exceed 400 mg/dL; a direct LDL measurement or the Martin-Hopkins equation should be used
  \item Non-fasting samples invalidate the triglyceride-derived LDL estimation and falsely alter triglyceride interpretation
  \item Acute illness, myocardial infarction, or surgery transiently lowers LDL-C; measure within 24 hours of an acute event or defer for 4--6 weeks
  \item Pregnancy elevates total cholesterol and triglycerides, so screening should be deferred during pregnancy
  \item Recent weight change or dietary change can alter results and should be noted
\end{itemize}
\section*{Risks}
Minimal. The only risks are those of routine venipuncture: transient pain, bruising, and, rarely, hematoma, infection, or vasovagal syncope.
\section*{What Happens After the Test}
Results are interpreted with the 10-year ASCVD risk estimate and current guideline thresholds to decide on lifestyle modification and statin intensity. Repeat testing is recommended 4--12 weeks after initiating or changing lipid-lowering therapy, and then periodically for monitoring. Triglycerides at or above 500 mg/dL warrant prompt treatment to reduce pancreatitis risk, and very high LDL-C prompts screening of first-degree relatives for familial hypercholesterolemia.
\section*{Understanding Results}
\subsection*{Desirable}
Total cholesterol <200 mg/dL. LDL-C <100 mg/dL (optimal <70--100 for high-risk/ASCVD; <55 mg/dL for very high-risk per 2018 guidelines) \cite{necp2001executive}. HDL-C $\geq$60 mg/dL (protective, subtract one risk factor). Triglycerides <150 mg/dL.
\subsection*{Borderline}
Total cholesterol 200--239 mg/dL. LDL-C 100--129 mg/dL (near optimal), 130--159 mg/dL (borderline high). Triglycerides 150--199 mg/dL.
\subsection*{High/Very High}
Total cholesterol $\geq$240 mg/dL. LDL-C 160--189 mg/dL (high), $\geq$190 mg/dL (very high). Triglycerides 200--499 mg/dL (high), $\geq$500 mg/dL (very high, risk of pancreatitis).
\section*{Normal Results}
Total cholesterol <200 mg/dL. LDL-C <100 mg/dL. HDL-C $\geq$40 mg/dL (men), $\geq$50 mg/dL (women). Non-HDL cholesterol (Total --- HDL) <130 mg/dL. Triglycerides <150 mg/dL. Fasting: 9--12 hours (water only). Non-fasting lipid panels (without TG calculation) are increasingly accepted for initial screening.
\section*{Follow-up Tests}
Apolipoprotein B (apoB) and lipoprotein(a) [Lp(a)], LDL particle number (NMR), high-sensitivity CRP (hs-CRP) for cardiovascular risk assessment, coronary artery calcium (CAC) score, and repeat lipid panel in 4--12 weeks after initiating or changing lipid-lowering therapy.
\section*{References}
\bibliographystyle{unsrtnat}
\bibliography{references}

\clearpage
\section*{Regulatory Status and Authorization Date}
Regulatory status applies to the specific assay, instrument, manufacturer, and intended use rather than to the test name alone. No single approval date can be assigned to this test category. For a commercial assay, verify the current product in the relevant regulator's database: the U.S. Food and Drug Administration (FDA) clearance, approval, or authorization; India's Central Drugs Standard Control Organisation (CDSCO) licence or permission; the European Union In Vitro Diagnostic Regulation (EU IVDR) CE certificate issued through the applicable conformity-assessment route; or the United Kingdom Medicines and Healthcare products Regulatory Agency (MHRA) registration and UKCA/CE status. Laboratory-developed tests may not have an individual FDA, CDSCO, EU IVDR, or MHRA approval date.
\clearpage
